% Options for packages loaded elsewhere
\PassOptionsToPackage{unicode}{hyperref}
\PassOptionsToPackage{hyphens}{url}
\documentclass[
  ignorenonframetext,
]{beamer}
\newif\ifbibliography
\usepackage{pgfpages}
\setbeamertemplate{caption}[numbered]
\setbeamertemplate{caption label separator}{: }
\setbeamercolor{caption name}{fg=normal text.fg}
\beamertemplatenavigationsymbolsempty
% remove section numbering
\setbeamertemplate{part page}{
  \centering
  \begin{beamercolorbox}[sep=16pt,center]{part title}
    \usebeamerfont{part title}\insertpart\par
  \end{beamercolorbox}
}
\setbeamertemplate{section page}{
  \centering
  \begin{beamercolorbox}[sep=12pt,center]{section title}
    \usebeamerfont{section title}\insertsection\par
  \end{beamercolorbox}
}
\setbeamertemplate{subsection page}{
  \centering
  \begin{beamercolorbox}[sep=8pt,center]{subsection title}
    \usebeamerfont{subsection title}\insertsubsection\par
  \end{beamercolorbox}
}
% Prevent slide breaks in the middle of a paragraph
\widowpenalties 1 10000
\raggedbottom
\AtBeginPart{
  \frame{\partpage}
}
\AtBeginSection{
  \ifbibliography
  \else
    \frame{\sectionpage}
  \fi
}
\AtBeginSubsection{
  \frame{\subsectionpage}
}
\usepackage{iftex}
\ifPDFTeX
  \usepackage[T1]{fontenc}
  \usepackage[utf8]{inputenc}
  \usepackage{textcomp} % provide euro and other symbols
\else % if luatex or xetex
  \usepackage{unicode-math} % this also loads fontspec
  \defaultfontfeatures{Scale=MatchLowercase}
  \defaultfontfeatures[\rmfamily]{Ligatures=TeX,Scale=1}
\fi
\usepackage{lmodern}
\ifPDFTeX\else
  % xetex/luatex font selection
\fi
% Use upquote if available, for straight quotes in verbatim environments
\IfFileExists{upquote.sty}{\usepackage{upquote}}{}
\IfFileExists{microtype.sty}{% use microtype if available
  \usepackage[]{microtype}
  \UseMicrotypeSet[protrusion]{basicmath} % disable protrusion for tt fonts
}{}
\makeatletter
\@ifundefined{KOMAClassName}{% if non-KOMA class
  \IfFileExists{parskip.sty}{%
    \usepackage{parskip}
  }{% else
    \setlength{\parindent}{0pt}
    \setlength{\parskip}{6pt plus 2pt minus 1pt}}
}{% if KOMA class
  \KOMAoptions{parskip=half}}
\makeatother
\setlength{\emergencystretch}{3em} % prevent overfull lines
\providecommand{\tightlist}{%
  \setlength{\itemsep}{0pt}\setlength{\parskip}{0pt}}
\usepackage{bookmark}
\IfFileExists{xurl.sty}{\usepackage{xurl}}{} % add URL line breaks if available
\urlstyle{same}
\hypersetup{
  hidelinks,
  pdfcreator={LaTeX via pandoc}}

\author{\texorpdfstring{}{}}
\date{}

\begin{document}

\begin{frame}[fragile]{Conclusion}
\protect\phantomsection\label{conclusion}
COGANT frames codebase analysis as a pipeline from ingestion through a
program graph IR to \textbf{Generalized Notation Notation (GNN)}
exports, with explicit confidence and provenance so that learning
systems can treat analysis noise as data rather than as hidden failure
modes. The Python layer provides session and pipeline APIs, a bundle
abstraction, CLI, review tooling, and HTML reporting; the Rust layer
concentrates graph mechanics and export formatting under crate
boundaries described in \texttt{../cogant/docs/ARCHITECTURE.md}.

\begin{block}{Shipped Capabilities}
\protect\phantomsection\label{shipped-capabilities}
Several capabilities ship in the current v0.1.x line and together define
the behaviour that downstream users can rely on:

\begin{enumerate}
\tightlist
\item
  \textbf{Fixpoint translation engine.} The shipped
  \texttt{TranslationEngine} re-applies every registered rule on each
  pass, terminates as soon as a pass produces zero new mappings, and
  bounds pathological rule sets with a configurable iteration cap. Each
  iteration boundary is recorded in the internal match log, so
  convergence can be audited after the fact.
\item
  \textbf{Rule priority and conflict resolution.} Overlapping mappings
  are reconciled by confidence score: when two mappings cover
  overlapping \texttt{graph\_fragment\_node\_ids}, the engine retains
  the higher-confidence mapping and records a
  \texttt{conflict\_resolved} event naming the winner, the loser, and
  the overlap. Rule priority is therefore expressed through the
  confidence model rather than a separate ordering table, and
  \texttt{translate\_with\_confidence()} re-scores and re-resolves so
  that priority shifts from the \texttt{ConfidenceModel} are honoured.
\item
  \textbf{Dynamic enrichment from coverage and traces.} When
  \texttt{.coverage} SQLite or Cobertura XML inputs and Chrome DevTools
  traces are supplied, \texttt{enrich\_graph()} annotates nodes with
  \texttt{coverage\_hits}, \texttt{branch\_coverage},
  \texttt{call\_count}, \texttt{avg\_duration\_ms}, and
  \texttt{is\_hot\_path} metadata, and appends
  \texttt{dynamic\_coverage} and \texttt{dynamic\_trace} markers to the
  program graph's evidence sources.
\item
  \textbf{Confidence tier promotion.} Mappings whose evidence set
  acquires dynamic markers become eligible for promotion from
  \texttt{STATIC\_ONLY} to \texttt{STATIC\_PLUS\_RUNTIME}, and hot-path
  and branch-coverage metadata raise the underlying score through the
  diversity bonus in Equation \ref{eq:confidence-core}. A purely-static
  run remains a first-class bundle; it is simply marked as such on the
  graph metadata.
\item
  \textbf{Ten-stage pipeline architecture.} The runner executes an
  ordered DAG (ingest, static, normalize, graph, translate, state space,
  process, export, validate, and report), with
  \texttt{skip\_on\_error=True} on optional stages so that a missing
  dynamic input or a pathological rule set produces a partial but
  well-formed Generalized Notation Notation bundle rather than an opaque
  failure.
\end{enumerate}

\textbf{Limitations} follow the honest scope in
\texttt{../cogant/docs/SPEC.md}: multi-language support beyond Python is
largely roadmap; translation rules and state-space extraction are
\textbf{partial} and repository-dependent; native acceleration is
staged. Users should validate exports on their own corpora before
trusting downstream model metrics.

\textbf{Intended users} include researchers building datasets from
open-source repositories, teams prototyping Active Inference models or
graph neural network training pipelines over program-graph data who need
a single export contract, and engineers extending the system via
\texttt{../cogant/docs/PLUGIN\_API.md}.

\textbf{Validation} in the software-engineering sense is split: the
repository's verification report enumerates implemented modules and
entry points; scientific validation of model quality remains the
responsibility of downstream training and evaluation code.
\end{block}

\begin{block}{Roadmap and Future Extensions}
\protect\phantomsection\label{roadmap-and-future-extensions}
Several concrete directions extend the current system:

\begin{enumerate}
\item
  \textbf{Multi-language parsers.} The v0.1.x front end targets Python;
  adding parsers for JavaScript/TypeScript, Java, and Go would cover the
  majority of open-source ML-relevant repositories. Each parser
  implements the plugin interface documented in
  \texttt{../cogant/docs/PLUGIN\_API.md}, so language additions do not
  require changes to the core IR or export pipeline.
\item
  \textbf{Rust acceleration of critical paths.} The native crate layer
  (\texttt{cogant-graph}, \texttt{cogant-translate}) already defines
  typed graph operations. Wiring PyO3 bindings for the hot paths ---
  deduplication, rule matching, and Generalized Notation Notation
  section/tensor packing in \texttt{cogant-gnn} --- would reduce
  end-to-end latency for large repositories from minutes to seconds,
  based on preliminary profiling of the Python fallback implementations.
\item
  \textbf{Incremental re-analysis.} Currently the pipeline processes a
  full repository snapshot on each invocation. An incremental mode that
  accepts a Git diff and updates only affected subgraphs would enable
  integration into CI/CD workflows where per-commit turnaround matters.
  The stable-identifier scheme already supports cross-run node matching;
  the missing piece is a dependency tracker that invalidates downstream
  IRs when upstream nodes change.
\item
  \textbf{LLM-assisted rule discovery.} Translation rules are currently
  hand-authored. A semi-automated workflow could present a large
  language model with unannotated graph fragments and ask it to propose
  candidate rules, which a human reviewer then accepts, edits, or
  rejects via the existing \texttt{ReviewAPI}. This combines the
  pattern-recognition strength of LLMs with the auditability of
  declarative rules.
\item
  \textbf{Cross-repository graph linking.} When multiple repositories
  share interfaces (for example, a library and its consumers), linking
  their program graphs at call boundaries produces a richer training
  signal for tasks such as API misuse detection and cross-project code
  search. The graph homomorphism property defined in Section 2 provides
  the formal basis for identifying shared interface nodes across
  independently analyzed repositories.
\end{enumerate}
\end{block}
\end{frame}

\end{document}
